\par\Needspace{6\baselineskip}\medskip
\SourceLeaf{224}\begin{center}\hypertarget{native-224-article161}{}\pdfbookmark[2]{161.}{article-161}161.\end{center}
We divide the \emph{first case} into three subcases.

I. If the point $B$ coincides with $A$ or with the opposite point, then we have $\delta=0$ or $=180^0$. The quantities $\gamma$, $\epsilon'$, $\epsilon''$ and the points $D'$, $D''$ will be indeterminate, while $\gamma'$, $\gamma''$, $\epsilon$ and the points $D$, $B^*$ will be determinate. The point $C$ necessarily coincides with $A$. By reasoning analogous to that in Article 140 we readily obtain the following equation.
\[
0=n'\frac{\sin(z-\sigma)}{\sin z}\cdot\frac{R'\sin\delta'}{R''\sin\delta''}
\cdot\frac{\sin(A''D-\delta'')}{\sin(A'D-\delta'+\sigma)}-n''.
\]
Everything set out in Articles 141, 142 may therefore be carried over to this case, provided we set $a=0$ and determine $b$ by equation 12 itself in Article 140. The quantities $z$, $r'$, $\dfrac{n'r'}{n}\mathpunct{\raisebox{.45ex}{,}}$
$\dfrac{n'r'}{n''}$ will be computed just as above. Once $z$ and hence the position of the point $C'$ are known, we can determine the position of the great circle $CC'$, its intersection with the great circle $A''B''$, namely the point $C''$, and therefore the arcs $CC'$, $CC''$, $C'C''$ or $2f''$, $2f'$, $2f$. This finally gives
\[
r=\frac{n'r'}{n}\cdot\frac{\sin2f}{\sin2f'}\mathpunct{\raisebox{.45ex}{,}}\qquad
r''=\frac{n'r'}{n''}\cdot\frac{\sin2f''}{\sin2f'}\mathpunct{\cdot}
\]

II. If the point $B''$ coincides with $A''$ or with the opposite point, everything we have just set out may be carried over, provided we interchange everything relating to the first position with what relates to the third.

III. The case in which $B'$ coincides either with $A'$ or with the opposite point must be treated somewhat differently. Here the point $C'$ will coincide with $A'$. The quantities $\gamma'$, $\epsilon$, $\epsilon''$ and the points $D$, $D''$, $B^*$ will be indeterminate. We can, however, determine the intersection of the great circle $BB''$ with the ecliptic\footnote{More generally, with the great circle $AA''$. For brevity, we consider here only the case in which the ecliptic is taken as the fundamental plane.}, and we set its longitude $=l'+\pi$. By reasoning similar to that developed in Article 140 we obtain the equation
\[
0=n\frac{R\sin\delta\sin(A''D'-\delta'')}{R''\sin\delta''\sin(AD'-\delta)}
+n'r'\frac{\sin\pi}{R''\sin(l''-l'-\pi)}+n''.
\]
Let us denote the coefficient of $n$, which is the same as $a$ in Article 140, by the same \SourceLeaf{225}symbol $a$, and the coefficient of $n'r'$ by $\beta$. We may also determine $a$ itself here by the following formula.
\[
a=-\frac{R\sin(l'+\pi-l)}{R''\sin(l''-l'-\pi)}
\]
Thus we have $0=an+\beta n'r'+n''$. If we combine this equation with
\[
P=\frac{n''}{n}\mathpunct{\raisebox{.45ex}{,}}\qquad
Q=2\left(\frac{n+n''}{n'}-1\right)r'^3,
\]
we obtain
\[
\frac{\beta(P+1)}{P+a}r'^4+r'^3+\tfrac12Q=0,
\]
from which we can derive the distance $r'$ except when $\beta=0$. In that case all that would follow is $P=-a$. Even outside the case $\beta=0$ (which would bring us to the third case, to be considered in the next article), $\beta$ will always be a very small quantity, so $P$ must differ little from $-a$. It is therefore evident that the determination of the coefficient $\dfrac{\beta(P+1)}{P+a}$ becomes very uncertain, so $r'$ cannot be determined with any precision.

Furthermore, we have $\dfrac{n'r'}{n}=-\dfrac{P+a}{\beta}\mathpunct{\raisebox{.45ex}{,}}$
$\dfrac{n'r'}{n''}=-\dfrac{P+a}{\beta P}\mathpunct{:}$ then, much as in Article 143, we readily derive the equations
\[
\begin{aligned}
r\sin\zeta&=\frac{n'r'}{n}\cdot\frac{\sin\gamma''}{\sin\epsilon'}\sin(l''-l')\\[.5em]
r''\sin\zeta''&=-\frac{n'r'}{n''}\cdot\frac{\sin\gamma}{\sin\epsilon'}\sin(l'-l)\\[.5em]
r\sin(\zeta-AD')&=r''P\frac{\sin\gamma''}{\sin\gamma}\sin(\zeta''-A''D'),
\end{aligned}
\]
Combining these with equations VIII and IX of Article 143 allows us to determine the quantities $r$, $\zeta$, $r''$, $\zeta''$. The remaining steps of the calculation agree with those described above.
